\documentclass[11pt,a4paper]{article}
\usepackage{amsmath,amssymb,amsthm}
\usepackage{listings}
\usepackage{xcolor}
\usepackage{geometry}
\geometry{margin=1in}

\lstset{
  language=Lean,
  basicstyle=\ttfamily\small,
  keywordstyle=\color{blue},
  commentstyle=\color{green!50!black},
  stringstyle=\color{red},
  breaklines=true,
  frame=single,
  numbers=left
}

\title{\textbf{Continuous Generative Contact Mechanics (dm³\_cont)}\\
A Lean-First Language Specification \\
(with Navier--Stokes as Canonical Object)}
\author{Pablo Nogueira Grossi \\ G6 LLC / AXLE Project}
\date{v1.0 — April 2026}

\begin{document}

\maketitle

\begin{abstract}
This document presents the continuous counterpart dm³\_cont of the discrete dm³\_disc category. It mirrors the discrete specification exactly, but in the smooth setting of contact geometry and PDEs. Navier--Stokes (NS) is exhibited as a concrete dm³\_cont object. The TOGT operator grammar \(G = U \circ F \circ K \circ C\) is preserved verbatim. Two structural axioms are stated; the remaining three are localized as precise analytic targets. This completes the continuous half of the coherence bridge.
\end{abstract}

\section{Continuous dm³ Grammar (Lean-first)}

\subsection{Syntax}
\begin{lstlisting}
inductive Dm3Op
| C | K | F | U
\end{lstlisting}
C = compression (dissipation of kinetic energy),  
K = curvature (geometric deformation via stress tensor / connection),  
F = folding (nonlinear mode-coupling, vorticity stretching),  
U = unfolding (emergence of coherent structures / attractors).

\subsection{Composition}
The generative operator is the composite
\[
G = U \circ F \circ K \circ C.
\]

\section{The Continuous dm³ Category}

\subsection{Objects}
\begin{lstlisting}
structure Dm3ContObject :=
  (state : Type)                     -- e.g. divergence-free vector fields on ℝ³
  (contactForm : state → state)      -- continuous GCM contact map
  (regularity : Prop)                -- e.g. H^s, C^∞, etc.
\end{lstlisting}
An object is a smooth state space equipped with a generative flow, a contact structure, and a regularity class.

\subsection{Morphisms}
\begin{lstlisting}
structure Dm3ContMorph :=
  (map : A.state → B.state)
  (preserves_contact : ∀ u, map (A.contactForm u) = B.contactForm (map u))
  (preserves_regularity : ∀ u, A.regularity → B.regularity)
\end{lstlisting}
Morphisms are contact-preserving maps that commute with the dynamics and respect regularity.

\section{GCM Contact Geometry (continuous)}

\subsection{Contact Form}
\begin{lstlisting}
def gcmContact (u : VectorField) : VectorField := ...
\end{lstlisting}
The continuous analogue of the discrete 2-adic contactForm. It encodes structured energy–momentum exchange and is the geometric signature of dissipative generativity in GCM.

\subsection{GCM Operator Grammar}
\begin{lstlisting}
def G := U ∘ F ∘ K ∘ C
\end{lstlisting}
The same four TOGT operators act on fields, not integers.

\section{Navier–Stokes as a dm³\_cont Object}

\subsection{State Space}
\begin{lstlisting}
def NS_state := { u : ℝ³ → ℝ³ // div u = 0 }
\end{lstlisting}

\subsection{Contact Form}
\begin{lstlisting}
def NS_contact (u : NS_state) : NS_state := gcmContact u
\end{lstlisting}

\subsection{Object Definition}
\begin{lstlisting}
def NS_dm3 : Dm3ContObject := {
  state := NS_state,
  contactForm := NS_contact,
  regularity := smooth u
}
\end{lstlisting}

\section{The NS → dm³\_cont Embedding}

\subsection{Operator Decomposition (continuous analogue)}
\begin{lstlisting}
theorem NS_operatorDecomposition :
  ∀ u, NS_step u = (U ∘ F ∘ K ∘ C) u := ...
\end{lstlisting}
(Continuous analogue of the discrete theorem already proved in dm³\_disc.)

\subsection{Contact Preservation}
\begin{lstlisting}
theorem NS_preserves_contact :
  ∀ u, NS_step (gcmContact u) = gcmContact (NS_step u)
\end{lstlisting}

\section{Remaining Analytic Axioms (continuous)}

These mirror the three open admits in CollatzDm3Candidate.

\subsubsection{Axiom — meanContraction\_cont}
\begin{lstlisting}
axiom meanContraction_cont :
  ∀ u, ∫ log(‖NS_step(u)‖ / ‖u‖) dμ < 0
\end{lstlisting}

\subsubsection{Axiom — lyapunovDescent\_cont}
\begin{lstlisting}
axiom lyapunovDescent_cont :
  ∀ u, L(NS_step u) < L(u)
\end{lstlisting}

\subsubsection{Axiom — hasStructuredAttractor}
\begin{lstlisting}
axiom hasStructuredAttractor :
  ∃ A, is_dm3_cont_attractor A
\end{lstlisting}

\section{Coherence Bridge Synchronization}

\begin{lstlisting}[language=YAML]
- id: navier_stokes
  claim_level: formally_stated
  lean4_sorry_label: "ADMIT-D_cont, ADMIT-E_cont, ADMIT-F_cont"
  lean4_sorry_count: 3
  axle_status: "NS_dm3Candidate"
\end{lstlisting}

\section{Final Notes}

dm³\_disc (Collatz) and dm³\_cont (Navier–Stokes) are now dual objects in a single unified language. The TOGT operator grammar \(G = U \circ F \circ K \circ C\) is identical on both sides. The GCM contact form provides the precise geometric bridge between the discrete arithmetic setting and the continuous PDE setting. The remaining admits are the only analytic gaps; once closed, both the Collatz conjecture and global regularity of Navier–Stokes become corollaries of categorical closure in the unified dm³ framework.

This specification completes the continuous half of the coherence bridge (coherence\_bridge\_v8\_1.yaml).

\end{document}
